\section{Scope, Chronology, and Evidence}

\subsection{Identity, names, and range}

Subject: Augustine of Hippo (Aurelius Augustinus), rhetorician, convert,
monk, presbyter, and bishop of Hippo Regius in Roman Numidia; Father and
Doctor of the Church; in received usage the Doctor of Grace. Whole-life
range: born at Thagaste (modern Souk Ahras, Algeria) 13 November 354; died
at Hippo Regius (modern Annaba, Algeria) 28 August 430, in the third month
of the Vandal siege. His mother's name, attested only once in his works, is
spelled \textit{Monnica} by nearly all manuscripts and \textit{Monica} in
received usage; this study writes Monica while recording the older
spelling. His companion of about fourteen years and the mother of his son
Adeodatus is unnamed in every ancient source. Modern claims about
Augustine's Berber, Punic, or Roman descent exceed the evidence and are
asserted nowhere in this study.

\subsection{Reader, question, and thesis boundary}

Written for a serious general reader, student, or researcher who wants one
coherent, evidence-controlled life. The governing questions are stated in
the opening section. The organizing synthesis---that the best-documented
conversion in ancient literature belongs to a witness who is also its
subject, and that the Church received both the doctor and his own published
self-corrections---is historical, not ecclesiastical. This study does not
adjudicate the theological merits of the Donatist, Pelagian, or
predestinarian questions; does not certify the garden scene's detail, the
relic chain, or any miracle; does not resolve the 395/396 consecration
year, the textual history of the Rule of St.~Augustine, or the exact
chronology of most works and letters; and claims no theological, canonical,
or rights review and no ecclesiastical approval.

\subsection{Corpus and citation conventions}

Augustine's \work{Confessions} are quoted from the public-domain Pilkington
translation (NPNF series 1, vol.~1, New York, 1886) in the New Advent
transcription, cited by book, chapter, and section (e.g., 8.12.29). His
letters follow the same volume's Cunningham translation with New Advent's
dates in parentheses; \work{On the Gift of Perseverance} follows NPNF
series 1, vol.~5. \work{The City of God} is quoted from the Dods
translation (Edinburgh, 1871) as retained in this repository's source
library, where the quoted passages are pinned to exact registered text
artifacts. Possidius's \work{Vita Augustini} is quoted from the
public-domain Weiskotten translation (Princeton, 1919) at tertullian.org,
cited by chapter. Paul the Deacon follows Foulke's translation (1907).
Vatican texts are quoted or paraphrased from the official English at
vatican.va. Quotations were checked against the online witnesses identified
in the References on 2026-07-25. Works are dated only where a stated basis
permits; no complete chronology of the sermons and letters is asserted.

\subsection{Evidence classes}

\begin{evidencekey}
\evidenceclass{A}{Subject-authored evidence: the \work{Confessions}
(prayerful retrospect of c.~397--400, not a contemporary record), letters,
treatises, and recorded acta; each claim limited by genre, audience, and
occasion.}
\evidenceclass{N}{Contemporary external witness: Possidius's \work{Vita}
(devoted, eyewitness for the clerical decades, dependent on Augustine's own
books for the early life); conciliar and imperial acts as reported in it.}
\evidenceclass{E}{Early received tradition: Paul the Deacon's
eighth-century relic notice; the early cult and transmission of the works.}
\evidenceclass{L}{Later tradition: the \work{Golden Legend} lives, the
seashore-child and \work{Te Deum} legends, iconography, later relic
recognitions.}
\evidenceclass{M}{Official ecclesial reception, limited to each act's
object: CCC 30; \work{Augustinum Hipponensem} (1986); Benedict XVI's Pavia
homily (2007) and general audiences (2008); the 1298 decretal as reported.}
\evidenceclass{K}{Named modern critical reconstruction: O'Donnell's
commentary (1992, online edition); the Stanford Encyclopedia entry; the
1907--1912 Catholic Encyclopedia; Brown 1967/2000 at publisher-record
level; each used to its stated premises with reporting ceilings marked.}
\evidenceclass{S}{Source-grounded project synthesis, never attributed to a
source.}
\end{evidencekey}

\subsection{Chronology}

Dates are claims with stated bases. Augustine's own narrative supplies ages
rather than years; the absolute skeleton rests on the birth and death dates
carried by official and critical tradition.

\begin{lifetimeline}
\lifeevent{13 Nov.\ 354}{Birth at Thagaste to Patricius and Monica}
{M (Benedict XVI, 9 Jan.\ 2008); A (ages in \work{Confessions})}{Secure}
\lifeevent{c.~365}{Boyhood illness; baptism requested and deferred}
{A (\work{Conf.} 1.11.17)}{Episode secure; date approximate}
\lifeevent{c.~370}{Idle sixteenth year at Thagaste; the pear theft}
{A (\work{Conf.} 2.3--2.4)}{Age stated by subject}
\lifeevent{c.~371}{Rhetoric at Carthage; death of Patricius (two years
before the \work{Hortensius})}{A (\work{Conf.} 3.1.1, 3.4.7)}{Approximate}
\lifeevent{c.~372}{Union with the unnamed concubine; birth of Adeodatus}
{A (\work{Conf.} 4.2.2; age ``barely fifteen'' at baptism, 9.6.14)}
{Inference from stated ages}
\lifeevent{373}{Reads the \work{Hortensius} in his nineteenth year; joins
the Manichees as hearer}{A (\work{Conf.} 3.4.7, 4.1.1)}{Secure relative
anchor}
\lifeevent{373--382}{Nine years a hearer; teaching at Thagaste and
Carthage; friend's death; first writings}{A (\work{Conf.} 4.1.1)}{Secure as
span}
\lifeevent{383}{Faustus at Carthage in his twenty-ninth year;
disillusion; sails for Rome}{A (\work{Conf.} 5.3.3, 5.7.13, 5.8.15)}
{Secure}
\lifeevent{384}{Milan chair of rhetoric through Symmachus; hears Ambrose}
{A (\work{Conf.} 5.13.23)}{Secure}
\lifeevent{385}{Monica reaches Milan; betrothal; concubine dismissed}
{A (\work{Conf.} 6.1, 6.13, 6.15)}{Approximate}
\lifeevent{386}{Platonist books; Victorinus story; Ponticianus; garden
resolution (late summer); withdrawal to Cassiciacum at the vintage
vacation}{A (\work{Conf.} 7--9); K (summer 386, SEP)}{Resolution secure;
scene single-witness}
\lifeevent{Easter 387}{Baptism at Milan by Ambrose, with Alypius and
Adeodatus (night of 23--24 April in the official statement)}{A
(\work{Conf.} 9.6.14); N (\work{Vita} 1); M (\work{Augustinum
Hipponensem})}{Secure}
\lifeevent{387 (after 25 Apr.)}{Ostia: the shared vision; death of Monica
aged 56}{A (\work{Conf.} 9.10--11); K (O'Donnell's bracket 25 Apr.--13
Nov.)}{Secure; day unknown}
\lifeevent{388}{Return to Africa (delayed by the Maximus war); community
at Thagaste ``for almost three years''}{N (\work{Vita} 2--3); K
(O'Donnell)}{Secure}
\lifeevent{c.~389/390}{Death of Adeodatus}{A (\work{Conf.} 9.6.14,
undated)}{Date unrecoverable}
\lifeevent{391}{Seized for ordination as presbyter at Hippo; study leave
requested; monastery founded}{N (\work{Vita} 4--5); A (\work{Letter} 21,
dated 391)}{Secure}
\lifeevent{392}{Public disputation with Fortunatus the Manichee}{N
(\work{Vita} 6; New Advent dating)}{Secure}
\lifeevent{395/396}{Consecration as coadjutor bishop by Megalius; sole
bishop on Valerius's death soon after}{N (\work{Vita} 8); M (395, Benedict
XVI); K (396 implied by CE's ``forty two'')}{Event secure; year
unresolved}
\lifeevent{c.~397--400}{\work{Confessions} written}{K (SEP ca.\ 396--400;
Benedict XVI 397--400)}{Approximate}
\lifeevent{399--c.~420}{\work{De Trinitate}: books 1--12 by 412, whole
revised c.~420}{M (Benedict XVI, 20 Feb.\ 2008)}{Approximate}
\lifeevent{405}{Imperial laws against heretics applied to the Donatists}
{K (CE ``Donatism'')}{Reported}
\lifeevent{408}{\work{Letter} 93 to Vincentius: the coercion change
documented}{A (New Advent date)}{Secure}
\lifeevent{410}{Sack of Rome by Alaric; refugees reach Africa; Pelagius
and Caelestius land near Hippo (411)}{M (Benedict XVI); K (CE
``Pelagianism'')}{Secure}
\lifeevent{411}{Conference of Carthage under Marcellinus; Donatists
adjudged and constrained}{N (\work{Vita} 13)}{Secure}
\lifeevent{412}{\work{Letter} 133: clemency plea for confessed
Circumcellion killers}{A (New Advent date)}{Secure}
\lifeevent{413--426}{\work{The City of God}, 22 books}{M (Benedict XVI,
20 Feb.\ 2008)}{Secure as span}
\lifeevent{415}{Synod of Diospolis accepts Pelagius's explanations}{K
(CE); A (\work{De gestis Pelagii}, cited at reported level)}{Secure}
\lifeevent{418}{Council of Carthage canons on sin and grace; Zosimus's
\work{Tractoria}}{N (\work{Vita} 18); K (CE)}{Secure}
\lifeevent{26 Sept.\ 426}{Eraclius designated successor; acta recorded}
{A (\work{Letter} 213); M (Benedict XVI, 16 Jan.\ 2008)}{Secure}
\lifeevent{c.~427}{\work{Retractationes} in two books}{M (Benedict XVI);
N (\work{Vita} 28)}{Approximate}
\lifeevent{429--430}{Vandal invasion of Africa; siege of Hippo begins
May/June 430}{N (\work{Vita} 28--29); M (Benedict XVI, 16 Jan.\ 2008)}
{Secure}
\lifeevent{28 Aug.\ 430}{Death in the third month of the siege;
``seventy-six years'' by inclusive count; burial at Hippo}{N (\work{Vita}
29, 31); M (Benedict XVI)}{Secure}
\lifeevent{c.~432+}{Possidius writes the \work{Vita} with its
\work{Indiculus}}{K (Weiskotten's dating, after July 431)}{Approximate}
\lifeevent{c.~6th c.?}{Relics carried to Sardinia ``on account of the
devastation of the barbarians''}{E (Paul the Deacon, undated); L (later
attribution to exiled bishops)}{Unresolved}
\lifeevent{by c.~720s}{Liutprand ransoms the relics from Sardinia; burial
at San Pietro in Ciel d'Oro, Pavia}{E (Paul the Deacon 6.48); L
(\work{Golden Legend} elaboration); M (Pavia homily, 2007)}{Secure as
8th-century report}
\lifeevent{1298}{Boniface VIII orders the feasts of the four Latin
doctors}{M/K (CE 1912, ``Doctors of the Church'')}{Secure}
\lifeevent{1430}{Monica's relics translated from Ostia to Rome under
Martin V}{K (CE, ``St.~Monica'')}{Reported}
\lifeevent{1483}{Caxton adds the seashore-child story to his \work{Golden
Legend} from an Antwerp altarpiece}{L (Caxton's own note)}{Secure as
textual fact}
\lifeevent{1986--2008}{\work{Augustinum Hipponensem}; papal visit to the
Pavia shrine (2007); five general audiences (2008)}{M (Vatican texts)}
{Secure}
\end{lifetimeline}

\subsection{Method, limits, and negative results}

Earliest controlled evidence governs; witnesses separated by centuries are
never silently merged; where the record conflicts or fails (the garden
scene's detail, 395/396, the Sardinia translation, the concubine's name),
the state of the question is printed beside the claim. Geographic and
linguistic limits: this study reads English translations of Latin sources;
no manuscript, epigraphic, or archaeological research was undertaken; no
Latin wording is asserted beyond famous tags identified as such; the
\work{Retractationes}, the Cassiciacum dialogues, \work{De gestis
Pelagii}, and the Divjak and Dolbeau discoveries are used only at the
reported level of the witnesses cited. Bounded negative results from
searches of the sources named in the References, conducted 2026-07-24/25:
no ancient source names Augustine's concubine; no source before the
eighth century records the movement of his relics; the seashore-child
story was not found in any source earlier than Caxton's 1483 note, and
its reported thirteenth-century exemplum carriers were not inspectable;
the \work{Te Deum} attribution was not found before the late
eighth-century tradition reported in the hymnological literature. These
are bounded, correctable search results, not proofs of absence. Mutable
online sources (vatican.va, New Advent, CCEL, tertullian.org,
elfinspell.com, Fordham sourcebooks, faculty.georgetown.edu, plato.stanford.edu,
ucpress.edu) were used as checked on those dates.

\subsection{Rights and review state}

Project prose is original; quotations of ancient authors follow the
public-domain NPNF translations (1886--1888), the public-domain Dods
translation (1871) retained in the repository source library, Weiskotten's
public-domain Possidius (1919), Foulke's Paul the Deacon (1907), and
Caxton's \work{Golden Legend} (Temple Classics text). Vatican texts
(Catechism, apostolic letter, homily, audiences) are quoted in brief
attributed snippets or paraphrased with citation; the O'Donnell commentary
and Stanford Encyclopedia are quoted in brief attributed snippets under
scholarly quotation practice; Brown is cited only at publisher-record
level. This publication has received internal editorial checking only: no
independent specialist, theological, or ecclesiastical review, and no
rights opinion, is claimed.
